% ---------------------------------------------------------------------------
% preamble.tex -- packages, theorem environments, boxes and macros shared by
% all lectures.  Only packages that ship with the local TinyTeX are used.
% ---------------------------------------------------------------------------

% ---- fonts and encoding ---------------------------------------------------
\usepackage[T1]{fontenc}
\usepackage[utf8]{inputenc}
\usepackage{lmodern}                   % Latin Modern text, sans, typewriter and math
\linespread{1.04}
\setlength{\emergencystretch}{2em}
% Latin Modern has no italic small caps: fall back to upright small caps silently
\input{t1lmr.fd}   % declare the family now so that the shape below can be added
\DeclareFontShape{T1}{lmr}{m}{scit}{<->ssub * lmr/m/sc}{}

% ---- mathematics ----------------------------------------------------------
\usepackage{amsmath,amssymb,amsthm,mathtools,bm}

% ---- layout ---------------------------------------------------------------
\usepackage[a4paper,left=3.1cm,right=3.1cm,top=3cm,bottom=3.2cm]{geometry}
\usepackage{booktabs,array,enumitem,float}
\usepackage{xcolor}
\usepackage{titlesec}
\setlist{itemsep=2pt,topsep=4pt,parsep=0pt}

\definecolor{accent}{HTML}{1B4F8A}     % deep blue
\definecolor{warn}{HTML}{B5470B}       % burnt orange
\definecolor{claim}{HTML}{8E1B3A}      % wine red, for unverified claims
\definecolor{softgray}{HTML}{666666}

% ---- headings -------------------------------------------------------------
\renewcommand{\chaptername}{Lecture}
\titleformat{\chapter}[display]
  {\sffamily\bfseries}
  {\color{accent}\large\MakeUppercase{\chaptertitlename}\ \thechapter}
  {4pt}
  {\raggedright\Huge}
  [\vspace{5pt}{\color{accent}\titlerule[1.2pt]}]
\titlespacing*{\chapter}{0pt}{-14pt}{22pt}
\titleformat{\section}
  {\sffamily\Large\bfseries\color{accent}}{\thesection}{0.7em}{}
\titleformat{\subsection}
  {\sffamily\large\bfseries\color{accent}}{\thesubsection}{0.7em}{}
\titleformat{\subsubsection}[runin]
  {\sffamily\bfseries}{}{0pt}{}[.]
\titlespacing*{\section}{0pt}{3.2ex plus 1ex minus .3ex}{1.4ex plus .2ex}
\titlespacing*{\subsection}{0pt}{2.6ex plus .8ex minus .3ex}{1ex plus .2ex}
\titlespacing*{\subsubsection}{0pt}{1.6ex plus .5ex}{0.8em}
\setcounter{secnumdepth}{2}
\setcounter{tocdepth}{1}
\pagestyle{plain}

% ---- theorem environments (one shared counter, numbered within a lecture) -
\theoremstyle{plain}
\newtheorem{theorem}{Theorem}[chapter]
\newtheorem{lemma}[theorem]{Lemma}
\newtheorem{proposition}[theorem]{Proposition}
\newtheorem{corollary}[theorem]{Corollary}
\newtheorem{conjecture}[theorem]{Conjecture}
\theoremstyle{definition}
\newtheorem{definition}[theorem]{Definition}
\newtheorem{example}[theorem]{Example}
\theoremstyle{remark}
\newtheorem*{remark}{Remark}
\newtheorem*{observation}{Observation}

\theoremstyle{definition}
\newtheorem{exercise}{Exercise}[chapter]

% ---- boxes ----------------------------------------------------------------
\usepackage{tikz}
\usetikzlibrary{arrows.meta,positioning,calc,shapes.geometric,fit,backgrounds,decorations.pathreplacing}
\usepackage[skins,breakable]{tcolorbox}

\tcbset{
  notebox/.style={
    enhanced, breakable, frame hidden, boxrule=0pt, arc=1.5pt,
    left=10pt, right=8pt, top=6pt, bottom=6pt,
    fonttitle=\sffamily\bfseries\small, coltitle=black, attach boxed title to top left={yshift=-2mm,xshift=0mm},
    before skip=10pt, after skip=10pt
  }
}
\newtcolorbox{bigpicture}[2][]{notebox, #1, colback=accent!6, colframe=accent,
  borderline west={2.6pt}{0pt}{accent}, title={\color{accent}#2}, before upper={\small\noindent}}
\newtcolorbox{pitfall}[2][]{notebox, #1, colback=warn!7, colframe=warn,
  borderline west={2.6pt}{0pt}{warn}, title={\color{warn}Pitfall: #2}}
\newtcolorbox{statusbox}[1]{notebox, colback=claim!6, colframe=claim,
  borderline west={2.6pt}{0pt}{claim}, title={\color{claim}Status note: #1}}
\newtcolorbox{goalbox}[1]{notebox, colback=black!4, colframe=black!50,
  borderline west={2.6pt}{0pt}{black!50}, title={#1}}

% ---- bibliography ---------------------------------------------------------
\usepackage[round,authoryear]{natbib}
\setlength{\bibsep}{3pt}

% ---- hyperlinks (load last) -----------------------------------------------
\usepackage[colorlinks=true,linkcolor=accent,citecolor=accent!70!black,urlcolor=accent,
            pdfusetitle,bookmarksnumbered]{hyperref}

% ---------------------------------------------------------------------------
% macros
% ---------------------------------------------------------------------------
\newcommand{\NP}{\mathsf{NP}}
\newcommand{\Ptime}{\mathsf{P}}
\newcommand{\coNP}{\mathsf{coNP}}
\newcommand{\PCP}{\mathsf{PCP}}
\newcommand{\BPP}{\mathsf{BPP}}
\newcommand{\poly}{\operatorname{poly}}
\newcommand{\polylog}{\operatorname{polylog}}
\newcommand{\eps}{\varepsilon}
\newcommand{\N}{\mathbb{N}}
\newcommand{\R}{\mathbb{R}}
\newcommand{\Z}{\mathbb{Z}}
\newcommand{\F}{\mathbb{F}}
\newcommand{\bits}{\{0,1\}}
\DeclareMathOperator*{\E}{\mathbb{E}}
\DeclareMathOperator{\OPT}{OPT}
\DeclareMathOperator{\val}{val}
\DeclareMathOperator{\Acc}{Acc}
\newcommand{\Sig}{\Sigma}
\newcommand{\UG}{\textsc{UG}}
\newcommand{\LC}{\textsc{LabelCover}}
\newcommand{\sat}[1]{\textsc{#1}}
\newcommand{\ie}{i.e.\ }
\newcommand{\eg}{e.g.\ }
\newcommand{\defn}[1]{\emph{#1}}

% Gap_{c,s}-Problem
\newcommand{\Gap}[3]{\textsc{Gap}_{#1,#2}\text{-}\textsc{#3}}

% a "star" difficulty marker for exercises
\newcommand{\stars}[1]{\textsf{\small\color{softgray}[#1]}}
